\section*{Chapter \ref{ch:m1:option-contracts-and-exchanges} --- Option Contracts and the Options Exchanges}

\begin{solution}{exo:m1:option-contracts-and-exchanges:1}
The XYZ put expiring on 15 January 2027 with a strike of \$47.50. The call:
\texttt{XYZ\textvisiblespace\textvisiblespace\textvisiblespace270319C00112500}.
\end{solution}

\begin{solution}{exo:m1:option-contracts-and-exchanges:2}
At 56: $10 \times 100 \times (6 - 2.40) = +\$3\,600$. At 48: the calls expire
worthless, $-\$2\,400$. Break-even at $50 + 2.40 = 52.40$.
\end{solution}

\begin{solution}{exo:m1:option-contracts-and-exchanges:3}
$2 \times (11 \times 70 + 5 \times 30) = 1\,840$ series. On sixteen exchanges:
29\,440 two-sided quotes to refresh for one cent in one share, which is why
options market making is first a problem of message rates and of deciding
which quotes not to update.
\end{solution}

\begin{solution}{exo:m1:option-contracts-and-exchanges:4}
47.00: both short puts are assigned, $+5\,000$ shares. 49.00: only the 50
puts, $+1\,000$. 50.00: nothing is \omterm{def:m1:option-contracts-and-exchanges:series}{in the money}, 0. 50.01: the long calls are
exercised, $+2\,500$. 52.50: still $+2\,500$. 52.51: the 60 short calls are
assigned, $2\,500 - 6\,000 = -3\,500$.
\end{solution}

\begin{solution}{exo:m1:option-contracts-and-exchanges:5}
Yes. Exercising captures the 0.60 dividend and gives up 0.15 of time value;
not exercising keeps an option whose underlying will open 0.60 lower. The
writer hedged with long shares is assigned overnight: the shares are called
away at 40, the dividend goes to the new owner, and the writer has lost the
0.60 it was counting on unless the option's price already reflected early
exercise. Desks run an early-exercise report every day before each ex-date.
\end{solution}

\begin{solution}{exo:m1:option-contracts-and-exchanges:6}
The first: the option cannot be traded after Thursday, yet its value is set
by Friday's opening prices. The desk carries the overnight gap and the
dispersion of the opening prints with no ability to adjust the option, only
its hedge, and must unwind the hedge \emph{in the opening auctions}. The
second: the option trades until the close, so the risk is continuous, but its
gamma in the final minutes is extreme near the strike, and the hedge must be
unwound in the closing auction.
\end{solution}

\begin{solution}{exo:m1:option-contracts-and-exchanges:7}
$+2\,500$ shares with exercise by exception, 0 with the contrary instruction.
Declining is right when the one cent of intrinsic value is less than the cost
and risk of the shares received: commissions and fees on 2\,500 shares, and
above all the exposure from Friday's close to Monday's open on shares bought
at 50 that may open at 49.50. For a holder who is not hedged short, one cent
does not pay for a weekend.
\end{solution}

\begin{solution}{exo:m1:option-contracts-and-exchanges:8}
The writer has assumed that holders exercise according to the 16:00 close.
They have until the evening cut-off and will exercise if the share trades
above 52.50 after hours. The writer does not know its position until the
next morning; with a hedge sized for ``expired worthless'' it may be short
6\,000 shares into Monday. At 15:59 the calls could have been bought back for
a few cents: the premium of certainty.
\end{solution}

\begin{solution}{pb:m1:option-contracts-and-exchanges:1}
\textbf{1.} Only the long 50 calls, by 2.48. The 52.50 calls are out of the
money by 0.02; both put series are out of the money.
\textbf{2.} The 25 long calls are exercised; everything else expires.
\textbf{3.} $+2\,500$ shares, bought at 50.
\textbf{4.} $-1\,900 + 2\,500 = +600$ shares.
\textbf{5.} The hedge was for a book that no longer exists: sell 600 shares,
ideally in Friday's closing auction, by estimating at 15:55 what exercise
will deliver.
\textbf{6.} Until the exchanges' cut-off, 16:30 Chicago time, 17:30 in New
York, and earlier at most brokers.
\textbf{7.} Because the share is worth 53.20 where it can now be sold, and
the call gives it for 52.50: exercise is worth 0.70 a share whatever the
16:00 print said.
\textbf{8.} 4\,500 shares delivered at 52.50.
\textbf{9.} $-1\,900 + 2\,500 - 4\,500 = -3\,900$ shares.
\textbf{10.} $4\,500 \times (54.10 - 52.50) = \$7\,200$.
\textbf{11.} $60 \times 100 \times 0.04 = \$240$.
\textbf{12.} Thirty times less.
\textbf{13.} Because of exactly this: the option does not expire at 16:00
but at the cut-off, and its seller is paid for ninety minutes of news and for
the buyer's right to decide afterwards.
\textbf{14.} It hedges the wrong scenario half of the time: if the stock
falls after hours nobody exercises, and the desk is long 3\,000 unwanted
shares into a lower open. Against Monday's 54.10 it would have gained
\$4\,860 on the shares and still been short 1\,500. Buying the options back
removes the uncertainty; buying shares bets on its outcome.
\textbf{15.} The desk holds the right: it decides, with the same ninety
minutes of hindsight, and can decline.
\textbf{16.} A \omterm{def:m1:what-a-trading-firm-does:market-maker}{market maker} long options near the strike is long gamma: it
sells shares as the price rises above the strike and buys as it falls below,
which pushes the price back towards the strike; the effect is strongest when
\omterm{def:m1:futures-contracts-and-exchanges:oi}{open interest} at that strike is large relative to the day's volume.
\textbf{17.} No: nothing is delivered, exercise is automatic against one
published settlement value, and there is no decision after the fact. The
risk moves into the settlement value itself (exercise 6).
\textbf{18.} With the client, as a short position in shares acquired at the
strike over the weekend; if the client cannot cover the margin, with the
broker, and beyond the broker with the clearing house's guarantee.
\textbf{19.} \textbf{Short 3\,900 shares.}
\textbf{20.} The holder of the option, after the close, with information the
writer also has but cannot act upon.
\end{solution}

\begin{solution}{iq:m1:option-contracts-and-exchanges:1}
American options can be exercised on any day up to \omterm{def:m1:option-contracts-and-exchanges:option}{expiry}, European only at
\omterm{def:m1:option-contracts-and-exchanges:option}{expiry}. In the US, options on shares and funds are American and physically
settled; the main index options are European and cash settled. The
difference matters around dividends and for deep in-the-money puts, where
early exercise is optimal, and for the writer's \omterm{def:m1:option-contracts-and-exchanges:assignment}{assignment} risk.
\iqlookfor{the product mapping and one consequence.}
\end{solution}

\begin{solution}{iq:m1:option-contracts-and-exchanges:2}
All US listed options are issued and cleared by one clearing house, so a
series is the same contract wherever it trades: long on one exchange and
short on another net to zero. A future is cleared by its own exchange's
clearing house, and an identical contract elsewhere is a different
counterparty.
\iqlookfor{clearing, not trading, as the source of fungibility.}
\end{solution}

\begin{solution}{iq:m1:option-contracts-and-exchanges:3}
A call: only to capture a dividend, just before the ex-date, when the
dividend exceeds the remaining time value (roughly, the value of the put at
the same strike plus interest on the strike). Otherwise selling the call is
better than exercising it. A put: when it is deep enough \omterm{def:m1:option-contracts-and-exchanges:series}{in the money} that
interest earned on the strike, received now, exceeds the remaining time
value; more often when rates are high.
\iqlookfor{the comparison with time value in both cases.}
\end{solution}

\begin{solution}{iq:m1:option-contracts-and-exchanges:4}
Being short options that expire with the underlying at the strike: the
number assigned, hence Monday's share position, is unknown until the next
morning, and holders decide after the close with later information. Manage
it by buying back short options near the strike before the close, by
flattening expiring strikes during the last days, by estimating \omterm{def:m1:option-contracts-and-exchanges:assignment}{assignment}
probabilities for what remains, and by keeping capacity to trade after hours
and at Monday's open.
\iqlookfor{closing the position rather than hedging the unknown.}
\end{solution}

\begin{solution}{iq:m1:option-contracts-and-exchanges:5}
An immutable record: underlying identifier, \omterm{def:m1:option-contracts-and-exchanges:option}{expiry} date, right, strike as an
integer in the smallest unit, multiplier, exercise style, settlement type,
and the deliverable for adjusted contracts. Never a float strike; never the
display symbol as the key, since roots change after corporate actions. Map
exchange symbols and OSI strings to the record through reference data with
effective dates; give each record a compact integer identifier for the hot
path.
\iqlookfor{integer strike, adjusted contracts, effective dating.}
\end{solution}

\begin{solution}{iq:m1:option-contracts-and-exchanges:6}
The option stops trading on Thursday; its payoff depends on a quotation built
from each stock's opening print. Hold the delta in futures or a basket
overnight; the exposure that remains is gamma against the overnight move,
which cannot be re-hedged in the option. On Friday morning the hedge must
be converted into exactly the settlement value: sell (or buy) the basket in
the opening auctions, or use futures that settle on the same quotation, which
removes the basis entirely. Size the position on Thursday with the overnight
gap, not the intraday volatility, in mind.
\iqlookfor{matching the hedge's exit to the settlement mechanism.}
\end{solution}
